\documentclass[12pt]{article}

\usepackage[a4paper,margin=2.5cm]{geometry}

\usepackage[onehalfspacing]{setspace}

\usepackage{amsmath}
\usepackage{amssymb}
\usepackage{xcolor}
\usepackage{mathtools}
\usepackage{amsthm}
\usepackage{tikz-cd}
\newtheorem{definition}{Definition}
\newtheorem{example}{Example}

\newenvironment{solution}
{\begin{proof}[Solution]}
{\end{proof}}

\newenvironment{Remark}
{\begin{proof}[Remark]}
{\end{proof}}

\setlength{\parindent}{0pt}

\begin{document}

\begin{titlepage}
    \centering

    \vspace*{\fill}

    {\Huge Linear Transformations II}

    \vspace{1cm}

    {\Large Jin-Ren Ryan Wang\par}

    \vspace{0.5cm}

    {\large July 31, 2026\par}

    \vspace*{\fill}
\end{titlepage}

\tableofcontents

\newpage

\section{Isomorphism of vector spaces}

\subsection{Definition}

\textbf{Definition 5.1.1. (inverse function)}
A function $f:A\to B$ is called \emph{bijective} if $f$ is both injective and surjective.
If $f:A\to B$ is bijective, then its inverse function $f^{-1}:B\to A$ exists and satisfies
$f^{-1}(f(a))=a,\quad \forall a\in A$ and $f\!\left(f^{-1}(b)\right)=b,\quad \forall b\in B$.

\medskip

\textbf{Definition 5.1.2.}
Let $V$ and $W$ be vector spaces over a field $F$.

\begin{description}
    \item[1.] $f:V\to W$ is called a \emph{(linear) isomorphism from $V$ onto $W$} if $f$ is a bijective\\
    linear map.

    \item[2.] We say that $V$ and $W$ are \emph{isomorphic as vector spaces} if there exists a linear\\
    isomorphism $f:V\to W$.
\end{description}

\subsection{Theorem}

\textbf{Theorem 5.1.1.}
Let $f:V\to W$ be a linear map of vector spaces. If $f$ is bijective 

(and hence $f^{-1}$ exists), then $f^{-1}:W\to V$ is also linear.

\begin{proof}
Since $f$ is surjective, every $\mathbf{w}\in W$ can be written as $\mathbf{w}=f(\mathbf{v})$ for some $\mathbf{v}\in V$.

We verify the two properties of linearity.

\begin{enumerate}
    \item \textbf{Additivity.}

    Let $\textcolor{red}{\mathbf{w}_1}=\textcolor{red}{f(\mathbf{v}_1)},\textcolor{red}{\mathbf{w}_2}=\textcolor{red}{f(\mathbf{v}_2)}$, where $\textcolor{red}{\mathbf{v}_1,\mathbf{v}_2\in V}$. Then
    \begin{align*}
    f^{-1}(\textcolor{red}{\mathbf{w}_1+\mathbf{w}_2})
        &=f^{-1}\!\left(\textcolor{red}{f(\mathbf{v}_1)+f(\mathbf{v}_2)}\right)\\
        &=f^{-1}\!\left(\textcolor{red}{f(\mathbf{v}_1+\mathbf{v}_2)}\right)\\
        &=\textcolor{red}{\mathbf{v}_1+\mathbf{v}_2}\\
        &=\textcolor{red}{f^{-1}(\mathbf{w}_1)+f^{-1}(\mathbf{w}_2)}.
    \end{align*}
    \item \textbf{Homogeneity.}

    Let $\lambda\in F$, $\textcolor{red}{\mathbf{w} \in W}$, \textcolor{red}{so $\mathbf{w}=f(\mathbf{v})$ for some $\mathbf{v}\in V$}. Then
    \begin{align*}
    f^{-1}(\lambda\textcolor{red}{\mathbf{w}})
        &=f^{-1}(\lambda \textcolor{red}{f(\mathbf{v})})\\
        &=f^{-1}(f(\lambda\textcolor{red}{\mathbf{v}}))\\
        &=\lambda\textcolor{red}{\mathbf{v}}\\
        &=\lambda \textcolor{red}{f^{-1}(\mathbf{w})}.
    \end{align*}
\end{enumerate}

This shows $f^{-1}:W\to V$ is also linear.
\end{proof}

\textbf{Theorem 5.1.2.}
Let $V$ be a vector space over a field $F$ of dimension $n$. Let $S=\{\mathbf{v}_1,\mathbf{v}_2,\ldots,\mathbf{v}_n\}$
be a basis of $V$. Then the function $\phi_S:V\to F^n$ defined by
\[
\mathbf{v}\mapsto[\mathbf{v}]_S.
\]
is a linear isomorphism.

\begin{proof}
Write
\[
\mathbf{v}=\sum_{i=1}^n\textcolor{red}{\alpha_i}\mathbf{v}_i,\qquad\mathbf{v}'=\sum_{i=1}^n\textcolor{red}{\alpha_i'}\mathbf{v}_i.
\]

We prove the statement in three steps.

\begin{enumerate}
    \item \fbox{$\phi_S$ is linear.}

    By definition, $[\mathbf{v}]_S=(\alpha_1,\alpha_2,\ldots,\alpha_n)$.

    Hence
    \begin{align*}
    \phi_S(\mathbf{v}+\mathbf{v}')
        &=\phi_S\!\left(
        \sum_{i=1}^n\textcolor{red}{(\alpha_i+\alpha_i')}\mathbf{v}_i
        \right)                                                   \\
        &=(\textcolor{red}{\alpha_1+\alpha_1',\ldots,\alpha_n+\alpha_n'})   \\
        &=(\alpha_1,\ldots,\alpha_n)+(\alpha_1',\ldots,\alpha_n') \\
        &=\phi_S(\mathbf{v})+\phi_S(\mathbf{v}').
    \end{align*}

    Likewise, for any $\lambda\in F$,
    \begin{align*}
    \phi_S(\lambda\mathbf{v})
        &=\phi_S\!\left(
        \sum_{i=1}^n(\lambda\textcolor{red}{\alpha_i})\mathbf{v}_i
        \right)                                     \\
        &=(\lambda\textcolor{red}{\alpha_1},\ldots,\lambda\textcolor{red}{\alpha_n}) \\
        &=\lambda(\textcolor{red}{\alpha_1,\ldots,\alpha_n})   \\
        &=\lambda\phi_S(\mathbf{v}).
    \end{align*}

    Therefore, $\phi_S$ is linear.

    \item \fbox{$\phi_S$ is injective.}

    Let $\mathbf{v}\in\ker(\phi_S)$ $\implies$ $\phi_S(\mathbf{v})=\textcolor{red}{(0,\ldots,0)}$.
    $\implies\mathbf{v} =\sum_{i=1}^n\textcolor{red}{0}\,\mathbf{v}_i =\mathbf{0}_V$.

    Thus
    \[
    \ker(\phi_S)=\{\mathbf{0}\},
    \]
    and therefore $\phi_S$ is injective.

    \item \fbox{$\phi_S$ is surjective.}

    By the Rank--Nullity Theorem,
    \[
    \operatorname{rank}(\phi_S)+\operatorname{nullity}(\phi_S)=\underbrace{\dim_F(V)}_{\textcolor{red}{\text{n}}}.
    \]

    Since $\phi_S$ is injective,
    \[
    \operatorname{nullity}(\phi_S)=0.
    \]

    Hence
    \[
    \operatorname{rank}(\phi_S)=\textcolor{red}{n}=\dim_F(V).
    \]

    Because $\dim_F(F^n)=n$, we conclude that $\operatorname{im}(\phi_S)=F^n$.
    Therefore, $\phi_S$ is surjective.
\end{enumerate}

Hence $\phi_S$ is a linear isomorphism.
\end{proof}

\textbf{Theorem 5.1.3.}
Let $V$ and $W$ be two finite-dimensional vector spaces over a field $F$. Then$V$ and $W$ are isomorphic as vector spaces
if and only if $\dim_F(V)=\dim_F(W)$.

\begin{proof}
We prove both directions.

\begin{description}
    \item[\textcolor{red}{($\Rightarrow$)}]
    Suppose $V$ and $W$ are isomorphic. By definition, $\exists$ a linear isomorphism $f:V\to W$

    Since $f$ is injective $\Longleftrightarrow$ $\operatorname{nullity}(f)=0$, and $f$ is surjective $\Longleftrightarrow$ $\operatorname{rank}(f)=\dim_F(W)$.

    By the Rank--Nullity Theorem, $\operatorname{rank}(f)+\operatorname{nullity}(f)=\dim_F(\textcolor{red}{V})$.

    Hence $\dim_F(W)+0=\dim_F(V)$, so $\dim_F(V)=\dim_F(W)$.

    \item[\textcolor{red}{($\Leftarrow$)}]
    Suppose $\dim_F(V)=\dim_F(W)=n$. Fix a basis $S$ of $V$, and a basis $S'$ of $W$.
    
    By Theorem 5.1.2, the coordinate maps
    \[
    \phi_{\textcolor{red}{S}}:V\to F^{\textcolor{red}{n}},\qquad\phi_{\textcolor{red}{S'}}:W\to F^{\textcolor{red}{n}}
    \]
    are linear isomorphisms.

    Now, $\begin{tikzcd} V\arrow[r, "\phi_{\textcolor{red}{S}}"]&F^{\textcolor{red}{n}}\arrow[r, "\phi_{\textcolor{red}{S'}}^{-1}"]&W \end{tikzcd}$

    Therefore, their composition $\phi_{S'}^{-1}\circ\phi_S:V\to W$ is also a linear isomorphism. 
    
    Hence, $V\cong W$. 

\end{description}
\end{proof}

\subsection{Example}

\subsubsection{Trivial Example}

Let $V$ be a vector space and let $\operatorname{id}_V:V\to V$ be the identity map on $V$. Then

\begin{solution}
Recall
\[
\operatorname{id}_V(\mathbf{v})=\mathbf{v},\quad \forall\,\mathbf{v}\in V.
\]

So

\begin{enumerate}
    \item $\operatorname{id}_V$ is a $\textcolor{red}{\text{linear map}}$.

    \item $\ker(\operatorname{id}_V)=\{\mathbf{0_V}\}$, hence $\operatorname{id}_V$ is $\textcolor{red}{\text{injective}}$.

    \item $\operatorname{im}(\operatorname{id}_V)=\underbrace{V}_{\textcolor{red}{\text{entire domain}}}$, hence $\operatorname{id}_V$ is $\textcolor{red}{\text{surjective}}$.
\end{enumerate}

Therefore, $\operatorname{id}_V:V\to V$ is a $\textcolor{red}{\text{linear isomorphism}}$. Consequently, V is isomorphic to itself.

Also, \fbox{\text{Every vector space is isomorphic to itself}}

\end{solution}

\section{Representations of Linear Maps as Matrices}

In this section, we describe a systematic way to represent every linear map 

$f:V\to W$ by a concrete matrix.

To do this, we begin a linear map $f:V\to W$ where $\dim_F(V)=n$ and $\dim_F(W)=m$.

Suppose we have fixed\
\[
B_1=\{\mathbf{v}_1,\mathbf{v}_2,\ldots,\mathbf{v}_n\} \text{ a basis of $V$}
\]
\[
B_2=\{\mathbf{w}_1,\mathbf{w}_2,\ldots,\mathbf{w}_m\} \text{ a basis of $W$}
\]

\medskip

Since $f$ is linear, it is completely determined by its values on the basis vectors of $V$. For each $\mathbf{v}_j\in B_1$, we can write
\[
f(\mathbf{v}_j)=\alpha_{1j}\mathbf{w}_1+\alpha_{2j}\mathbf{w}_2+\cdots+\alpha_{mj}\mathbf{w}_m,\qquad j=1,\ldots,n.
\]

That is, we express each $f(\mathbf{v}_j)$ as a linear combination of the basis vectors of $W$.

\medskip

We now collect the coefficients into a matrix. Specifically, we define
\[
[f]_{B_1}^{B_2}=(\alpha_{ij})_{i,j}=
\begin{pmatrix}
\alpha_{11} & \alpha_{12} & \cdots & \alpha_{1n}\\
\alpha_{21} & \alpha_{22} & \cdots & \alpha_{2n}\\
\vdots      & \vdots      & \ddots & \vdots\\
\alpha_{m1} & \alpha_{m2} & \cdots & \alpha_{mn}
\end{pmatrix},
\]
whose $j$-th column is exactly the coordinate vector of $f(\mathbf{v}_j)$ with respect to the basis $B_2$.

\subsection{Definition}

\textbf{Definition 5.2.1.}
The matrix constructed above is called the \emph{matrix representing
$f$ with respect to the bases $B_1$ and $B_2$}. We denote it by $[f]_{B_1}^{B_2}$.

\textbf{Definition 5.2.2.}
Let $A = (\alpha_{ij})$ be an $l\times m$ matrix and $B = \beta_{ij}$ be an $m \times n$ matrix.

Then the product AB is a $l \times n$ matrix $C = (\gamma_{ij})$ where,
\[
(i,j)\text{-entry of }AB = \gamma_{ij} = \sum_{\textcolor{red}{k}=1}^{m} \alpha_{i\textcolor{red}{k}}\beta_{\textcolor{red}{k}j}
\]

\subsection{Theorem}

\textbf{Theorem 5.2.1.}
Let $f:V\to W$ be a linear map between finite-dimensional vector spaces over a field $F$.
Fix a basis $B_1$ be a basis of $V$, and $B_2$ be a basis of $W$.

Then, for every $\mathbf{v}\in V$, we have the equality of matrices (and vectors):
\[
[f]_{B_1}^{B_2}[\mathbf{v}]_{B_1}=[f(\mathbf{v})]_{B_2}.
\]

\begin{proof}
Let $B_1=\{\mathbf{v}_1,\ldots,\mathbf{v}_n\}$, $B_2=\{\mathbf{w}_1,\ldots,\mathbf{w}_m\}$

Let
\[
\textcolor{red}{[f]_{B_1}^{B_2}=(\alpha_{ij})_{i,j} \Longleftrightarrow f(\mathbf{v}_j)=\sum_{i=1}^{m}\alpha_{ij}\mathbf{w}_i,\forall j}
\]

And let
\[
[\mathbf{v}]_{B_1}=
\begin{pmatrix}
\lambda_1\\
\vdots\\
\lambda_n
\end{pmatrix} \Longleftrightarrow \mathbf{v}=\sum_{j=1}^{n}\lambda_j\mathbf{v}_j
\]

Then
\begin{align*}
[f]_{B_1}^{B_2}[\mathbf{v}]_{B_1}
=
(\alpha_{ij})_{i,j}
\begin{pmatrix}
\lambda_1\\
\vdots\\
\lambda_n
\end{pmatrix} 
=
\begin{pmatrix}
\sum_{j=1}^{n}\alpha_{1j}\lambda_j\\
\vdots\\
\sum_{j=1}^{n}\alpha_{mj}\lambda_j
\end{pmatrix}.
\end{align*}
On the other hand, since $f$ is linear,
\begin{align*}
f(\mathbf{v})
&=
f\left(
\sum_{j=1}^{n}\lambda_j\textcolor{red}{\mathbf{v}_j}
\right)                                              
=
\sum_{j=1}^{n}\lambda_j\textcolor{red}{f(\mathbf{v}_j)}            
=
\sum_{j=1}^{n}\lambda_j
\left(
\textcolor{red}{\sum_{i=1}^{m}\alpha_{ij}\mathbf{w}_i}
\right)                              
=
\textcolor{red}{\sum_{i=1}^{m}}
\left(
\sum_{j=1}^{n}\alpha_{ij}\lambda_j
\right)
\textcolor{red}{\mathbf{w}_i}.
\end{align*}

Hence,
\[
[f(\mathbf{v})]_{B_2}=
\begin{pmatrix}
\sum_{j=1}^{n}\alpha_{1j}\lambda_j\\
\vdots\\
\sum_{j=1}^{n}\alpha_{mj}\lambda_j
\end{pmatrix},
\]
which is exactly the vector obtained above. Therefore,
\[
[f]_{B_1}^{B_2}[\mathbf{v}]_{B_1}=[f(\mathbf{v})]_{B_2}.
\]

\begin{Remark}
The theorem can be summarized by the following commutative diagram:
\[
\begin{array}{ccc}
V & \xrightarrow{\;f\;} & W\\
\downarrow{\scriptstyle\phi_{B_1}} && \downarrow{\scriptstyle\phi_{B_2}}\\
F^n & \textcolor{red}{\xrightarrow{\;[f]_{B_1}^{B_2}\;}} & F^m
\end{array}
\]

That is,
\[
\textcolor{red}{[f]_{B_1}^{B_2}}\circ\phi_{B_1}=\phi_{B_2}\circ f,
\]
or equivalently,
\[
\textcolor{red}{[f]_{B_1}^{B_2}}[\mathbf{v}]_{B_1}=[f(\mathbf{v})]_{B_2},\qquad\forall\,\mathbf{v}\in V.
\]
\end{Remark}

\end{proof}

\subsection{Example}

\subsubsection{Matrix Representation of Differentiation}

Consider the linear map $D:\mathbb{R}[x]_{\le 3}\to \mathbb{R}[x]_{\le 2}$ defined by $p(x)\mapsto p'(x)$.
Fix the bases 
\[
S_1=\{1,x,x^2,x^3\}\text{ of }\mathbb{R}[x]_{\le 3},\text{ and } S_2=\{1,x,x^2\}\text{ of }\mathbb{R}[x]_{\le 2}.
\]
Find the matrix $[D]_{S_1}^{S_2}$

\begin{solution}
We compute the images of the basis vectors of $S_1$ 

and express them in terms of the basis $S_2$.

\medskip

First,
\[
D(\textcolor{red}{1})=0=\textcolor{red}{0}\cdot1+\textcolor{red}{0}\cdot x+\textcolor{red}{0}\cdot x^2,
\]
so
\[
[D(1)]_{S_2}=(0,0,0)^T.
\]

Next,
\[
D(\textcolor{red}{x})=1=\textcolor{red}{1}\cdot1+\textcolor{red}{0}\cdot x+\textcolor{red}{0}\cdot x^2,
\]
so
\[
[D(x)]_{S_2}=(1,0,0)^T.
\]

Similarly,
\[
D(\textcolor{red}{x^2})=2x=\textcolor{red}{0}\cdot1+\textcolor{red}{2}\cdot x+\textcolor{red}{0}\cdot x^2,
\quad
[D(x^2)]_{S_2}=(0,2,0)^T,
\]

and
\[
D(\textcolor{red}{x^3})=3x^2=\textcolor{red}{0}\cdot1+\textcolor{red}{0}\cdot x+\textcolor{red}{3}\cdot x^2,
\quad
[D(x^3)]_{S_2}=(0,0,3)^T.
\]

\medskip

Placing these coordinate vectors as columns, we obtain
\[
[D]_{\textcolor{red}{S_1}}^{\textcolor{red}{S_2}}
=
\begin{pmatrix}
0 & 1 & 0 & 0\\
0 & 0 & 2 & 0\\
0 & 0 & 0 & 3
\end{pmatrix}.
\]

\begin{Remark}
For example, if
\[
p(x)=\textcolor{red}{1+2x+4x^3},\qquad D[p(x)]=\textcolor{red}{2 + 12x^2}
\]
then
\[
[p(x)]_{S_1}=
\begin{pmatrix}
\textcolor{red}{1}\\
\textcolor{red}{2}\\
\textcolor{red}{0}\\
\textcolor{red}{4}
\end{pmatrix},
\qquad
[D(p(x))]_{S_2}=
\begin{pmatrix}
\textcolor{red}{2}\\
\textcolor{red}{0}\\
\textcolor{red}{12}
\end{pmatrix}.
\]

We observe that
\[
[D(p(x))]_{S_2}=[D]_{S_1}^{S_2}[p(x)]_{S_1}.
\]

Thus differentiation has been translated into matrix multiplication.
\end{Remark}

\end{solution}

\subsubsection{Matrix Representations under Different Bases}

Consider the linear map $f:\mathbb{R}^3\rightarrow\mathbb{R}^3$, defined by
\[
(x,y,z)\longmapsto
(-x-2y+2z,\,
4x+3y-4z,\,
-2y+z).
\]

\begin{enumerate}
    \item[(a)] Verify that $f$ is a linear transformation.
    \item[(b)] Find the matrix representing $f$ with respect to the
    standard basis of $\mathbb{R}^3$ (for both the domain and codomain).
    \item[(c)] Verify that
    \[
    S=\{(1,0,1),(0,1,1),(1,-1,1)\}
    \]
    is a basis of $\mathbb{R}^3$. Find the matrix $[f]_S^S$.
\end{enumerate}

\newpage

\begin{solution}
The following are the solutions of three subproblems

\medskip

\text{(a)}

For any $\mathbf{u}=(x_1,y_1,z_1)$, $\mathbf{v}=(x_2,y_2,z_2)\in\mathbb{R}^3$,
\[
\begin{aligned}
f(\mathbf{u}+\mathbf{v})
&=
f(x_1+x_2,\,
y_1+y_2,\,
z_1+z_2) \\
&=
\Bigl(
-(x_1+x_2)-2(y_1+y_2)+2(z_1+z_2),\\
&\qquad
4(x_1+x_2)+3(y_1+y_2)-4(z_1+z_2),\\
&\qquad
-2(y_1+y_2)+(z_1+z_2)
\Bigr) \\
&=
f(\mathbf{u})+f(\mathbf{v}).
\end{aligned}
\]

Moreover, for any $\lambda\in\mathbb{R}$,
\[
\begin{aligned}
f(\lambda\mathbf{u})
&=
f(\lambda x_1,\lambda y_1,\lambda z_1)   \\
&=
\Bigl(
-\lambda x_1-2\lambda y_1+2\lambda z_1,\,
4\lambda x_1+3\lambda y_1-4\lambda z_1,\,
-2\lambda y_1+\lambda z_1
\Bigr)    \\
&=
\lambda f(\mathbf{u}).
\end{aligned}
\]

Therefore, $f$ satisfies both additivity and homogeneity. Hence, $f:\mathbb{R}^3\to\mathbb{R}^3$
is a linear transformation.

\medskip

\text{(b)}
Let $B=\{\mathbf e_1,\mathbf e_2,\mathbf e_3\}$ be the standard basis of $\mathbb R^3$.

Compute the images of the basis vectors:
\[
f(\textcolor{red}{\mathbf e_1})=(-1,4,0)=\textcolor{red}{(-1)}\quad \mathbf e_1+\textcolor{red}{4}\quad \mathbf e_2+\textcolor{red}{0}\quad \mathbf e_3,
\]
\[
f(\textcolor{red}{\mathbf e_2})=(-2,3,-2)=\textcolor{red}{(-2)}\quad \mathbf e_1+\textcolor{red}{3}\quad \mathbf e_2+\textcolor{red}{(-2)}\quad \mathbf e_3,
\]
\[
f(\textcolor{red}{\mathbf e_3})=(2,-4,1)=\textcolor{red}{1}\quad \mathbf e_1+\textcolor{red}{-4}\quad \mathbf e_2+\textcolor{red}{1}\quad \mathbf e_3.
\]
Hence,
\[
[f]_B^B=
\begin{pmatrix}
-1&-2&2\\
4&3&-4\\
0&-2&1
\end{pmatrix}.
\]

\medskip

\text{(c)}

Let
\[
S=\{\mathbf v_1,\mathbf v_2,\mathbf v_3\}=\{(1,0,1),(0,1,1),(1,-1,1)\}.
\]

Compute
\[
\begin{aligned}
f(\textcolor{red}{\mathbf v_1})
    &=(1,0,1)=\textcolor{red}{1}\mathbf v_1+\textcolor{red}{0}\mathbf v_2+\textcolor{red}{0}\mathbf v_3,\\
f(\textcolor{red}{\mathbf v_2})
    &=(0,-1,-1)=\textcolor{red}{0}\mathbf v_1+\textcolor{red}{(-1)}\mathbf v_2+\textcolor{red}{0}\mathbf v_3,\\
f(\mathbf v_3)
    &=(3,-3,3)=\textcolor{red}{0}\mathbf v_1+\textcolor{red}{0}\mathbf v_2+\textcolor{red}{3}\mathbf v_3.
\end{aligned}
\]

Therefore,
\[
[f]_S^S=
\begin{pmatrix}
1&0&0\\
0&-1&0\\
0&0&3
\end{pmatrix}.
\]

\begin{Remark}
Although the \textcolor{red}{moral standard basis is convenient} for computation, 

\textcolor{red}{it is not necessarily the "nicest" basis for representing a linear transformation}.

Question: How to find the "nice" basis for a linear map f?

Answer: Diagonalization Theory
\end{Remark}

\end{solution}

\subsubsection{Using Matrix Representations to Study a Linear Map}

Consider the function $f:\mathbb{R}[x]_{\le2}\longrightarrow\mathbb{R}[x]_{\le3}$
\[
p(x)\longmapsto-\frac{d}{dx}\bigl(p(x^2-1)\bigr)+\int_2^x p(t)\,dt.
\]

\begin{enumerate}
    \item[(a)] Verify that $f$ is a linear transformation.
    \item[(b)] Find $\ker(f)$. Is $f$ injective?
    \item[(c)] Find $\operatorname{rank}(f)$. Is $f$ surjective?
\end{enumerate}

\begin{solution}{
The following are the solutions of abobe three subproblems

\medskip

\text{(a)}

Let $p(x),q(x)\in\mathbb{R}[x]_{\le2}$ and let $\lambda\in\mathbb{R}$.

Observe that differentiation, composition with $x^2-1$, and integration
are all linear operators. Hence

\begin{align*}
f(p+q)
&=
-\frac{d}{dx}\!\left((p+q)(x^2-1)\right)
+\int_2^x(p(t)+q(t))\,dt      \\
&=
-\frac{d}{dx}\bigl(p(x^2-1)\bigr)
-\frac{d}{dx}\bigl(q(x^2-1)\bigr)\\
&\qquad
+\int_2^xp(t)\,dt
+\int_2^xq(t)\,dt             \\
&=
f(p)+f(q).
\end{align*}

Similarly,
\begin{align*}
f(\lambda p)=-\frac{d}{dx}\!\left(\lambda p(x^2-1)\right)+\int_2^x\lambda p(t)\,dt
=
\lambda
\left(
-\frac{d}{dx}(p(x^2-1))
+\int_2^xp(t)\,dt
\right)
=
\lambda f(p).
\end{align*}

Therefore, $f$ satisfies additivity and homogeneity. Hence,
$f$ is a linear transformation.
}

\medskip

\text{(b)}

Fix a basis $S_1=\{1,x,x^2\}$ \textcolor{red}{(domain)}, and a basis $S_2=\{1,x,x^2,x^3\}$ \textcolor{red}{(co-domain)}.

Compute the images of the basis vectors:
\[
f(\textcolor{red}{1})=x-2 = \textcolor{red}{(-2)} + \textcolor{red}{1}x + \textcolor{red}{0}x^2 + \textcolor{red}{0}x^3
\]
\[
f(\textcolor{red}{x})=-2x+\frac12x^2-2 =\textcolor{red}{(-2)} + \textcolor{red}{(-2)}x + \textcolor{red}{\frac12}x^2 + \textcolor{red}{(\frac{-11}{3})}x^3
\]
\[
f(\textcolor{red}{x^2})=-\frac13x^3+4x-\frac83 = \textcolor{red}{(\frac{-8}{3})} + \textcolor{red}{4}x + \textcolor{red}{0}x^2 + \textcolor{red}{(\frac{-11}{3})}x^3
\]

Hence

\[
A \coloneq [f]_{S_1}^{S_2}=
\begin{pmatrix}
-2&-2&\frac{-8}3\\
1&-2&4\\
0&\frac12&0\\
0&0&\frac{-11}3
\end{pmatrix}.
\]

Since

\[
\ker(f)
\cong
\ker([f]_{S_1}^{S_2}),
\]

and the above matrix has trivial nullspace,

\[
\ker(A)=\ker([f]_{S_1}^{S_2})=\{\mathbf0\}=
\left\{
\begin{pmatrix}
0\\
0\\
0\\
\end{pmatrix}
\right\}
\]

we conclude that

\[
\ker(f)=\{0\textcolor{red}{\text{ (zero polynomial)}}\},
\]

where $0$ denotes the zero polynomial.

Therefore, $f$ is injective.

\text{(c)}

By the Rank--Nullity Theorem,
\[
\operatorname{rank}(f)+\operatorname{nullity}(f)=\textcolor{red}{\dim\!\left(\mathbb{R}[x]_{\le2}\right)}.
\]

Since $\ker(f)=\{0\}$, we have $\operatorname{nullity}(f)=0$.
$\implies\operatorname{rank}(f)=\textcolor{red}{3}$, however, $\dim\!\left(\mathbb{R}[x]_{\le3}\right)=\textcolor{red}{4}$

$\implies\operatorname{rank}(f)\neq\dim\!\left(\textcolor{red}{\underbrace{\mathbb{R}[x]_{\le3}}_{\text{co-domain}}}\right)$

$\implies$ $f$ is NOT surjective.
\end{solution}

\newpage

\section{Composition of linear maps and matrix multiplication}

\subsection{Definition}

\textbf{Definition 5.3.1.}
Let $f:A\to B$ and $g:B\to C$ be two functions.
The \emph{composition} of $f$ and $g$ is the function $g\circ f:A\to C$ is defined by
\[
(g\circ f)(a)=g(f(a)),\qquad \forall\,a\in A.
\]

For $f:A\to A$ then we denote $f^n\coloneq\underbrace{f\circ f\circ\cdots\circ f}_{n\text{ times}}$.

\subsection{Theorem}

\textbf{Theorem 5.3.1.}
Let $f:V\to U$ and $g:U\to W$ be linear maps between finite-dimensional vector spaces over a field
$F$. Fix
\[
\text{a basis $B_1$ of $V$ },\text{a basis $B_2$ of $U$ },\text{a basis $B_3$ of $W$ }
\]

Then we have an equality of matrices:
\[
[g\circ f]_{B_1}^{B_3}=[g]_{B_2}^{B_3} \cdot [f]_{B_1}^{B_2}.
\]

\begin{proof}

Let
\begin{align*}
B_1=\{\mathbf{v}_1,\mathbf{v}_2,\ldots,\mathbf{v}_m\} \text{ be a basis of $V$}\\
B_2=\{\mathbf{u}_1,\mathbf{u}_2,\ldots,\mathbf{u}_n\} \text{ be a basis of $U$}\\
B_3=\{\mathbf{w}_1,\mathbf{w}_2,\ldots,\mathbf{w}_p\} \text{ be a basis of $W$}
\end{align*}

Let
\begin{align*}
[f]_{B_1}^{B_2}=(\alpha_{ij})_{i,j} \text{ which means } f(\mathbf{v}_j)=\sum_{k=1}^{n}\alpha_{kj}\mathbf{u}_k\\
[g]_{B_2}^{B_3}=(\beta_{ij})_{i,j} \text{ which means } g(\mathbf{u}_k)=\sum_{i=1}^{p}\beta_{ik}\mathbf{w}_i
\end{align*}

We now determine the matrix
\[
[g\circ f]_{B_1}^{B_3}.
\]

For each basis vector $\mathbf{v}_j$,

\begin{align*}
(g\circ f)(\textcolor{red}{\mathbf{v}_j})
&=
g\!\left(f(\mathbf{v}_j)\right)                                           \\
&=
g\!\left(
\sum_{k=1}^{n}\alpha_{kj}\textcolor{red}{\mathbf{u}_k}
\right)                                                                    \\
&=
\sum_{k=1}^{n}\alpha_{kj}\textcolor{red}{g(\mathbf{u}_k)}
\qquad
(\text{since }g\text{ is linear})                                          \\
&=
\sum_{k=1}^{n}\alpha_{kj}\textcolor{red}{\sum_{i=1}^{p}\beta_{ik}\mathbf{w}_i}                                      \\
&=
\textcolor{red}{\sum_{i=1}^{p}}
\left(
\sum_{k=1}^{n}\beta_{ik}\alpha_{kj}
\right)
\textcolor{red}{\mathbf{w}_i}.
\end{align*}

Hence, the $(i,j)$-entry of $[g\circ f]_{B_1}^{B_3}$ is $\sum_{k=1}^{n}\beta_{ik}\alpha_{kj}$,
which is exactly the $(i,j)$-entry of the matrix product $[g]_{B_2}^{B_3} \cdot[f]_{B_1}^{B_2}$.

Therefore,
\[
[g\circ f]_{B_1}^{B_3}=[g]_{B_2}^{B_3}\cdot[f]_{B_1}^{B_2}.
\]

\end{proof}

\subsection{Example}

\subsubsection{Composition of Matrix Multiplication}

The point $(a,b)$ in $\mathbb{R}^2$ is first rotatedcounterclockwise by $\frac{\pi}4$ about the origin and then
reflected about the line $y=\sqrt3\,x$. Find the resulting point.


\begin{solution}

Let 
\begin{align*}
\text{$B$ be the standard basis of $\mathbb{R}^2$}\\
\textcolor{red}{T_1}:\text{rotation by }\frac{\pi}{4}\text{ (anti-clockwise)}\\
T_2:\text{reflection along the line }y=\textcolor{red}{\tan\!\left(\frac{\pi}{3}\right)}x
\end{align*}

Then
\[
[\textcolor{red}{T_1}]_B^B=
\begin{pmatrix}
\cos\frac{\pi}{4} & -\sin\frac{\pi}{4}\\
\sin\frac{\pi}{4} & \cos\frac{\pi}{4}
\end{pmatrix}
=
\begin{pmatrix}
\dfrac1{\sqrt2} & -\dfrac1{\sqrt2}\\[2mm]
\dfrac1{\sqrt2} & \dfrac1{\sqrt2}
\end{pmatrix}.
\]
\[
[T_2]_B^B=
\begin{pmatrix}
\cos\frac{2\pi}{3} & \sin\frac{2\pi}{3}\\
\sin\frac{2\pi}{3} & -\cos\frac{2\pi}{3}
\end{pmatrix}
=
\begin{pmatrix}
-\dfrac12 & \dfrac{\sqrt3}{2}\\[2mm]
\dfrac{\sqrt3}{2} & \dfrac12
\end{pmatrix}.
\]

\newpage

Hence, 
\[
[T_2\circ \textcolor{red}{T_1}]_B^B=[T_2]_B^B\cdot[\textcolor{red}{T_1}]_B^B \textcolor{red}{\text{ (matrix multiplication)}}
\]
\[
=
\begin{pmatrix}
-\dfrac12 & \dfrac{\sqrt3}{2}\\[2mm]
\dfrac{\sqrt3}{2} & \dfrac12
\end{pmatrix}
\begin{pmatrix}
\dfrac1{\sqrt2} & -\dfrac1{\sqrt2}\\[2mm]
\dfrac1{\sqrt2} & \dfrac1{\sqrt2}
\end{pmatrix}
\]
\[
=
\begin{pmatrix}
\dfrac{-1+\sqrt3}{2\sqrt2}
&
\dfrac{1+\sqrt3}{2\sqrt2}
\\[3mm]
\dfrac{1+\sqrt3}{2\sqrt2}
&
\dfrac{1-\sqrt3}{2\sqrt2}
\end{pmatrix}.
\]

Therefore,
\[
\begin{pmatrix}
a\\
b
\end{pmatrix}
\longmapsto
\begin{pmatrix}
\dfrac{-1+\sqrt3}{2\sqrt2}
&
\dfrac{1+\sqrt3}{2\sqrt2}
\\[3mm]
\dfrac{1+\sqrt3}{2\sqrt2}
&
\dfrac{1-\sqrt3}{2\sqrt2}
\end{pmatrix}
\begin{pmatrix}
a\\
b
\end{pmatrix}
=
\underbrace{
\begin{pmatrix}
\dfrac{-1+\sqrt3}{2\sqrt2}a
+\dfrac{1+\sqrt3}{2\sqrt2}b
\\[3mm]
\dfrac{1+\sqrt3}{2\sqrt2}a
+\dfrac{1-\sqrt3}{2\sqrt2}b
\end{pmatrix}
}_{\textcolor{red}{\text{in standard basis of B}}}.
\]

\end{solution}

\newpage

\section{Inverse of a matrix}

\subsection{Definition}

\textbf{Definition 5.4.1.}
Let $A\in M_n(F)$ be a square matrix. If there exists a matrix $B\in M_n(F)$ s.t. $AB=BA=I_n$,
then
\begin{enumerate}
    \item we say that $A$ is \emph{invertible} (or \emph{non-singular}); and

    \item we call $B$ the \emph{inverse} of $A$ and denote it by $A^{-1}$.
\end{enumerate}

\subsection{Theorem}

\textbf{Theorem 5.4.1.}
Let $A\in M_{m\times n}(F)$. Suppose there exists $B\in M_{n\times m}(F)$ s.t.
\[
BA=I_{\textcolor{red}{n}}\text{ and }AB=I_{\textcolor{red}{m}}.
\]

Then $n=m=\operatorname{rank}(A)$.

\begin{proof}

We shall show that $\operatorname{rank}(A)=m=n$.

\text{(1)} First observe that
\[
\operatorname{rank}(A)=\operatorname{rank}(A\overbrace{\textcolor{red}{BA}}^{\textcolor{red}{I_n}})\le\operatorname{rank}(AB)=\operatorname{rank}(\textcolor{red}{I_m})=m,
\]

On the other hand,
\[
m=\operatorname{rank}(I_m)=\operatorname{rank}(AB)\le\operatorname{rank}(A).
\]

Therefore,
\[
\operatorname{rank}(A)=m.
\]

\medskip

\text{(2)} Similarly,
\[
\operatorname{rank}(A)=\operatorname{rank}(ABA)\le\operatorname{rank}(BA)=\operatorname{rank}(\textcolor{red}{I_n})=n,
\]

Also,
\[
n=\operatorname{rank}(I_n)=\operatorname{rank}(BA)\le\operatorname{rank}(A).
\]

Hence,
\[
\operatorname{rank}(A)=n.
\]

Therefore, if $A$ has an inverse, then
\[
\textcolor{red}{n=m=\operatorname{rank}(A)}.
\]

\begin{Remark}
Rank inequalities
\[
\left\{   
\begin{aligned}     
\operatorname{rank}(f\circ g)\le\operatorname{rank}(f)\\
\operatorname{rank}(f\circ g)\le\operatorname{rank}(g)
\end{aligned}
\right.  
\]

and \textcolor{red}{matrix version}
\[
\textcolor{red}{
\left\{   
\begin{aligned}     
\operatorname{rank}(AB)\le\operatorname{rank}(A)\\
\operatorname{rank}(AB)\le\operatorname{rank}(B)
\end{aligned}
\right.  
}
\]
\end{Remark}

\end{proof}

\textbf{Theorem 5.4.2.}
Let $f:V\to W$ be an isomorphism of finite-dimensional vector spaces.
Fix any basis \(B_1\) of \(V\) and \(B_2\) of \(W\).
Then we have an equality of matrices
\[
\left([f]_{B_1}^{B_2}\right)^{-1}=[f^{-1}]_{\textcolor{red}{B_2}}^{\textcolor{red}{B_1}}.
\]

In other words, the inverse of a matrix represents the inverse of a linear map.

\begin{proof}

By Theorem 5.3.1,
\[
[f^{-1}]_{B_2}^{B_1}[f]_{B_1}^{B_2}=[f^{-1}\circ f]_{B_1}^{B_1}.
\]

Since $f^{-1}\circ f=\operatorname{id}_V$, we obtain
\[
[f^{-1}\circ f]_{B_1}^{B_1}=[\operatorname{id}_V]_{B_1}^{B_1}=I_n,\text{ where }n=\dim(V)
\]
Therefore,
\[
[f^{-1}]_{B_2}^{B_1}=\left([f]_{B_1}^{B_2}\right)^{-1}.
\]

\end{proof}

\section{Remark}
\begin{center}
\renewcommand{\arraystretch}{1.4}
\begin{tabular}{|c|c|}
\hline
\textbf{Vector space / Linear map world}
&
\textbf{Vector / Matrix world}
\\
\hline

$f:V\rightarrow W$
&
$[f]_{B_1}^{B_2}$ \textcolor{red}{ \text{(matrix)}}
\\
\hline

Composition
$f\circ g$
&
Multiplication $[f\circ g]_{\textcolor{red}{B_1}}^{\textcolor{red}{B_3}}=[f]_{\textcolor{red}{B_2}}^{\textcolor{red}{B_3}}[g]_{\textcolor{red}{B_1}}^{\textcolor{red}{B_2}}$
\\
\hline

Addition
$f+g$,
Scaling
$\lambda f$
&
Addition
$[f+g]_{\textcolor{red}{B_1}}^{\textcolor{red}{B_2}}=[f]_{\textcolor{red}{B_1}}^{\textcolor{red}{B_2}}+[g]_{\textcolor{red}{B_1}}^{\textcolor{red}{B_2}}$,
Scaling
$[\lambda f]_{\textcolor{red}{B_1}}^{\textcolor{red}{B_2}}=\lambda[f]_{\textcolor{red}{B_1}}^{\textcolor{red}{B_2}}$
\\
\hline

Inverse
$f^{-1}:W\rightarrow V$
&
Matrix inverse
$[f^{-1}]_{\textcolor{red}{B_2}}^{\textcolor{red}{B_1}}=([f]_{\textcolor{red}{B_1}}^{\textcolor{red}{B_2}})^{-1}$
\\
\hline
\end{tabular}

\end{center}
Question: When is a matrix invertible?

\end{document}